\subsubsection{Subproceso 1: Acreditación y Pre-requisitos}

\begin{enumerate}
    \item \textbf{Cierre de actas de notas y envío de convalidaciones previas.} \\
    \textit{Responsable: Facultad / Escuela Profesional y DIRCE.} \\
    Se cierran las calificaciones del semestre anterior sin reclamos pendientes y se procesan en el SIGA-ORCE los cambios de plan de estudios y convalidaciones aprobadas.

    \item \textbf{Pago del autoseguro estudiantil y regularización de deudas.} \\
    \textit{Responsable: Estudiante Regular / Unidad de Tesorería y Dependencias de Facultad.} \\
    El alumno cancela el aporte obligatorio del autoseguro médico y subsana cualquier deuda de libros o materiales en bibliotecas y laboratorios de su Facultad.

    \item \textbf{Registro de nóminas de tutoría y de matrícula preferencial.} \\
    \textit{Responsable: Jefatura de Tutoría y Dirección de Escuela Profesional.} \\
    Los tutores autorizan la activación del código de los estudiantes en riesgo académico que asistieron a consejería, y la Escuela Profesional remite el listado de alumnos egresantes con 30 créditos o menos por cursar.

    \item \textbf{Validación automática de pagos, adeudos y bloqueos.} \\
    \textit{Responsable: Sistema SIGA-ORCE / Intranet DIRCE.} \\
    La plataforma cruza la información de recaudación y de la Facultad para comprobar que el autoseguro figure pagado y que el estudiante no tenga bloqueos ni adeudos activos.

    \item \textbf{Cálculo del promedio ponderado, tope de créditos y orden de mérito.} \\
    \textit{Responsable: DIRCE / Sistema SIGA-ORCE.} \\
    El sistema calcula el promedio ponderado de los dos últimos semestres regulares del estudiante, le asigna su límite máximo de créditos y determina su posición en el orden de mérito.

    \item \textbf{Asignación de grupos y publicación de turnos de matrícula.} \\
    \textit{Responsable: Comisión de Matrícula y DIRCE.} \\
    Se organiza a los estudiantes en sus respectivos grupos de atención (Riesgo Académico, Promoción y Grupos Regulares por promedio) y se publican las fechas y horas de inicio y fin en el portal.

    \item \textbf{Generación de la ficha de matrícula y habilitación en la Intranet DIRCE.} \\
    \textit{Responsable: Intranet DIRCE / Sistema SIGA-ORCE.} \\
    El portal carga la ficha individual con el plan de estudios, especialidad y créditos máximos del alumno, dejando habilitado su acceso para cuando inicie su turno en el Subproceso 2.
\end{enumerate}
